\documentclass[12pt]{article}

\usepackage[margin=1in]{geometry}
\usepackage{booktabs}
\usepackage{array}
\usepackage{tabularx}
\usepackage{hyperref}
\usepackage{amsmath}

\title{Empirical Counterparts for MVPF Terms}
\author{Taco Prins, Lennard Schlattmann, and Daniel Jonas Schmidt}
\date{\today}

\begin{document}

\maketitle

\section*{Empirical Estimates by Term}

\begin{table}[h]
\centering
\small
\renewcommand{\arraystretch}{1}
\begin{tabularx}{\linewidth}{p{4.6cm}p{4.3cm}X}
\toprule
\textbf{Term} & \textbf{Estimate} & \textbf{Source} \\
\midrule
$I$ (take-up rate)
    & $\approx 30\%$ in high-risk areas
    & \href{https://www.fema.gov/openfema-data-page/nfip-residential-penetration-rates-v1}{FEMA OpenFEMA NFIP Penetration Rates}; \href{https://www.rand.org/content/dam/rand/pubs/technical_reports/2006/RAND_TR300.pdf}{Dixon et al.\ (2006)} \\
$\varepsilon$ (price semi-elasticity)
    & $-0.17$ to $-1.38$
    & \href{https://faculty.wharton.upenn.edu/wp-content/uploads/2017/07/MismeasuringRisk_Mulder2021.pdf}{Mulder (2021)};
      \href{https://journalofcrr.com/research/03-07-gourevitch-et-al/}{Gourevitch et al.\ (2023)};
      \href{https://link.springer.com/article/10.1023/A:1007823631497}{Browne \& Hoyt (2000)} \\
$\partial I / \partial a$ (crowd-out)
    & $-3\%$ of coverage
    & \href{https://www.sciencedirect.com/science/article/abs/pii/S0095069617303479}{Kousky, Michel-Kerjan \& Raschky (2018)} \\
$F(q)$ (belief distribution)
    & Substantial heterogeneity
    & \href{https://www.nber.org/system/files/working_papers/w23854/w23854.pdf}{Bakkensen \& Barrage (2017)};
      \href{https://www.aeaweb.org/articles?id=10.1257/app.6.3.206}{Gallagher (2014)};
      \href{https://faculty.wharton.upenn.edu/wp-content/uploads/2017/07/MismeasuringRisk_Mulder2021.pdf}{Mulder (2021)} \\
$p$ (true flood probability)
    & $1\%$--$10\%$ per year
    & \href{https://zenodo.org/records/6459076}{First Street Foundation (Zenodo v2.0)};
      \href{https://www.fema.gov/flood-maps/national-flood-hazard-layer}{FEMA National Flood Hazard Layer} \\
$d$ (flood damage, conditional on flood)
    & \$44k--\$83k (avg.\ claim)
    & \href{https://www.fema.gov/openfema-data-page/fima-nfip-redacted-claims-v2}{OpenFEMA NFIP Redacted Claims};
      \href{https://www.floodsmart.gov/historical-nfip-claims-information-and-trends}{NFIP Historical Claims (FloodSmart)} \\
$d/w$ (damage-to-wealth ratio)
    & $10\%$--$16\%$
    & \href{https://www.fema.gov/openfema-data-page/fima-nfip-redacted-claims-v2}{OpenFEMA NFIP Redacted Claims};
      \href{https://www.floodsmart.gov/historical-nfip-claims-information-and-trends}{NFIP Historical Claims (FloodSmart)} \\
$w$ (household wealth)
    & \$193k (median); \$300k--\$500k (home value)
    & \href{https://fortune.com/2023/10/18/us-household-wealth-increase-survey-of-consumer-finances-federal-reserve/}{Federal Reserve SCF (2022)};
      \href{https://www.census.gov/programs-surveys/ahs.html}{American Housing Survey} \\
$s$ (current subsidy rate)
    & $\approx 47\%$ (median policyholder)
    & \href{https://www.gao.gov/assets/gao-23-105977.pdf}{GAO (2023)};
      \href{https://www.fema.gov/flood-insurance/risk-rating}{FEMA Risk Rating 2.0} \\
$a$ (current relief fraction)
    & $\approx 4\%$--$7\%$
    & Computed from avg.\ FEMA IA grant (\$3k) over avg.\ claim (\$44k--\$83k);
      \href{https://www.sciencedirect.com/science/article/abs/pii/S0095069617303479}{Kousky et al.\ (2018)} \\
$\lambda$ (shadow value of budget)
    & $0.5$--$2$
    & \href{https://academic.oup.com/qje/article/135/3/1209/5781614}{Hendren \& Sprung-Keyser (2020)} \\
\bottomrule
\end{tabularx}
\caption{Empirical counterparts for the terms in MVPF$_s$ and MVPF$_a$. All estimates are from US data except where noted. The crowd-out estimate from Kousky et al.\ (2018) is identified on the intensive margin (coverage amount) rather than the binary take-up margin, partly due to a regulatory requirement that disaster aid recipients maintain insurance. The model calibration of $d/w = 0.5$ exceeds the empirical average of 10--16\%; the latter reflects average NFIP claims (\$44,000--\$83,000) relative to average US home values (\$500,000).}
\end{table}

\section*{Notes on Identification}

\paragraph{Price semi-elasticity $\varepsilon$.}
Since the household's premium is $\pi = (1-s)pd$, we have $\partial\pi/\partial s = -pd$ and therefore
\[
    \frac{\partial I}{\partial s} = -pd \cdot I \cdot \varepsilon,
\]
where $\varepsilon \equiv (1/I)(\partial I / \partial \pi) < 0$.
Mulder (2021) exploits variation from outdated FEMA flood maps that misclassify high-risk homes as low-risk, yielding $|\varepsilon| \approx 0.17$.
Gourevitch et al.\ (2023) exploit the NFIP's Risk Rating~2.0 reform and find a wider range of $|\varepsilon|$ from 0.15 to 1.38 across subgroups.

\paragraph{Crowd-out $\partial I / \partial a$.}
Direct quasi-experimental estimates of $\partial I/\partial a$ are rare because variation in disaster relief generosity is uncommon.
Kousky et al.\ (2018) link FEMA Individual Assistance grants (2000--2011) to NFIP policy data and find that receiving a grant reduces coverage the following year by \$4{,}000--\$5{,}000 ($\approx 3\%$ of mean coverage).
Alternatively, using the ratio trick (eq.~7 of the model notes),
\[
    \frac{\partial I/\partial a}{\partial I/\partial s} = -\frac{q^*}{p}\cdot\frac{u'(c_{U,F})}{u'(c_I)},
\]
$\partial I/\partial a$ can be recovered from the estimated $\varepsilon$ together with structural utility assumptions.

\paragraph{Belief distribution $F(q)$.}
Gallagher (2014) documents that take-up spikes after a flood and then decays, consistent with $q < p$ on average and with beliefs that revert toward prior underestimation over time.
Bakkensen and Barrage (2017) document large cross-sectional heterogeneity in subjective flood probabilities among US coastal homeowners.
Mulder (2021) measures subjective flood risk at the property level alongside objective risk, enabling direct estimation of the belief gap $p - q$.

\paragraph{Risk aversion $\gamma$.}
Chetty (2006) identifies $\gamma \approx 1$ from labor supply responses, providing a lower bound.
The broader empirical literature yields estimates in $[1, 4]$; we use $\gamma = 2$ which is in the middle of this range.

\paragraph{Shadow value $\lambda$.}
At the interior optimum, MVPF$_s$ = MVPF$_a$ = $\lambda$.
Hendren and Sprung-Keyser (2020) compute MVPFs for 133 US policy changes; their estimates for adult-targeted social insurance programs cluster between 0.5 and 2, providing a plausible external benchmark for $\lambda$.

\section*{Numerical MVPF Estimates at Current US Policy}

Using the empirical parameter values above, we evaluate both MVPFs at the current US policy mix ($s = 0.47$, $a = 0.055$, $d/w = 0.15$).
The crowd-out response $\partial I/\partial a$ is recovered via the ratio trick.
Table~\ref{tab:mvpf} reports all intermediate quantities and the resulting MVPFs.

\begin{table}[h]
\centering
\renewcommand{\arraystretch}{1.4}
\begin{tabular}{lc}
\toprule
& \textbf{Value} \\
\midrule
\multicolumn{2}{l}{\textit{Parameters}} \\
$p$           & $0.02$ \\
$d/w$         & $0.15$ \\
$s$           & $0.47$ \\
$a$           & $0.055$ \\
$I$           & $0.30$ \\
$\varepsilon$ & $-0.17$ \\
$\gamma$      & $2.0$ \\
\midrule
\multicolumn{2}{l}{\textit{Consumption levels}} \\
$c_I$     & $0.9984$ \\
$c_{U,F}$ & $0.8583$ \\
\midrule
\multicolumn{2}{l}{\textit{Intermediate quantities}} \\
$q^*$                         & $0.0096$ \\
$(p - q^*)\,\Delta u$         & $0.00171$ \\
$u'(c_I)$                     & $1.0032$ \\
$u'(c_{U,F})$                 & $1.3576$ \\
$\partial I/\partial s$       & $0.000153$ \\
$\partial I/\partial a$       & $-0.000100$ \\
\midrule
\multicolumn{2}{l}{\textit{MVPFs}} \\
$\mathrm{MVPF}_s$ & $\mathbf{1.003}$ \\
$\mathrm{MVPF}_a$ & $\mathbf{1.358}$ \\
\bottomrule
\end{tabular}
\caption{MVPF estimates at current US policy ($s = 0.47$, $a = 0.055$, $I = 0.30$, $\varepsilon = -0.17$, $\gamma = 2$, $p = 0.02$, $d/w = 0.15$).}
\label{tab:mvpf}
\end{table}

$\mathrm{MVPF}_a > \mathrm{MVPF}_s$, implying that the current US policy mix over-invests in insurance subsidies relative to disaster relief.
The subsidy MVPF is close to 1 because insured households retain nearly all their wealth ($c_I \approx w$), leaving little room for a welfare gain beyond the direct transfer.
The relief MVPF is substantially larger because flood victims suffer a significant consumption loss ($c_{U,F} \ll w$), so their marginal utility $u'(c_{U,F})$ is much higher.
The internality correction $(p - q^*)\Delta u$ and fiscal externality terms are negligibly small, so the MVPFs are essentially determined by the direct transfer terms $pd\,u'(c_I)\,I$ and $pd\,u'(c_{U,F})\,(1-I)$.

\paragraph{Utility function.}
All computations use CRRA utility with $\gamma = 2$:
\[
    u(c) = \frac{c^{1-\gamma}}{1-\gamma} = -\frac{1}{c}, \qquad u'(c) = c^{-\gamma} = \frac{1}{c^2}.
\]
Utility is defined over total final consumption, with wealth normalized to $w = 1$.
The consumption levels entering the MVPF formulas are therefore
\begin{align*}
    c_I     &= w - (1-s)\,pd = 1 - 0.53 \times 0.003 = 0.9984, \\
    c_{U,F} &= w - (1-a)\,d = 1 - 0.945 \times 0.15 = 0.8583.
\end{align*}
Because the actuarially fair premium $pd = 0.02 \times 0.15 = 0.003$ is small, insured households barely move away from $w = 1$, so $u'(c_I) \approx 1$.
This mechanically drives $\mathrm{MVPF}_s \approx 1$: the subsidy delivers one unit of welfare per unit of fiscal cost, with almost no amplification from the utility curvature.
The gap $\mathrm{MVPF}_a - \mathrm{MVPF}_s$ is therefore driven almost entirely by the consumption loss of flood victims, i.e.\ by $c_{U,F} = 0.8583 < 1$, which raises $u'(c_{U,F}) = 1.36 > 1 \approx u'(c_I)$.

\end{document}
